\section*{September 7, 2026: Final choices for the three groups}
\textit{Agent-drafted record of Jacob's confirmed decisions.}

Jacob finalized the 97,215-square-foot Assessor denominator for the combined
84-unit first phase at 4400 Grove. The commercial decision input now uses that
area, retains the two parcel identifiers, and identifies the 2021 Assessor as
its land source. Holding the existing floor area fixed implies floor-area ratio
1.1734 and 37.6386 units per acre. This replaces the prior 196,020-square-foot
choice in the ledger; it has not yet changed the paper's frozen dataset.

Jacob also accepted the later Assessor year 2006 for 1217 W Arthington. The review
input labels this an accepted proxy and preserves the older 2005 evidence.
Accepting the year does not resolve the original four-versus-six unit conflict;
the existing six-unit count remains unchanged.

Natchez A, B, and C will remain three separate construction groups, with 72,
84, and 39 units and years 2017, 2020, and 2020 respectively. All three already
exist separately in the frozen dataset. A and B pass the main 500-foot sample
restrictions at approximately 377 and 420 feet; C is outside at 657 feet.
Each observation contains multiple buildings: 12 in A, 14 in B, and 8 in C.
The obsolete six-building description for C was corrected, but its land area
and excluded parcel still need checking. No new split, merger, or exclusion
was introduced.

\small
\textit{Reproduction.} The three decision tables are now tracked at their
production paths. Comparing them with their research-archive versions confirms
that only these five reviewed rows changed and that row counts remain 80, 47,
and 179. The commercial build cannot reach the revised project ledger because
\path{commercial_city_building_footprints.gpkg} has no producer rule. The
logbook was rebuilt through Make with the same three unrelated audit outputs
held fixed. This is a recorded decision update, not a completed reconstruction
or new regression result.
\normalsize
